\subsection{RQ4: Can proprietary LLM detect SATD in test code?}
\textbf{Motivation.} Motivated by the unexpectedly poor performance of open-source LLMs observed in RQ3, we extended our investigation to proprietary LLMs. Previous research consistently reports that proprietary models tend to outperform open-source counterparts across a range of standard benchmarks, including those that involve code understanding and generation tasks~\cite{dvivedi2024comparative, mayer_can_2024, jimenez2023swe}.

\textbf{Approach.} We evaluated two major proprietary LLM families—Gemini and GPT—to assess their effectiveness in detecting SATD in test code. For Gemini, we used the latest Flash variants of Gemini 2.0 and 2.5, and for GPT, we tested three GPT-5 models (nano, mini, and standard) to capture diverse reasoning capabilities and model sizes. Each model was tested under the few-shot configurations of 0, 2, and 4 shots. All the experiments were conducted using the prompt shown in Figure~\ref{fig:satd-detection-prompt-template}. 
% We excluded keyword-based prompts, including the best-performing MAT-suggested prompts, for proprietary LLMs, as their use with open-source LLMs demonstrated low recall—particularly in failing to detect SATD comments that lacked explicit keyword indicators.


\begin{table*}[!h]
\centering
\caption{Precision (P), Recall (R), F1-score (F1) and Matthews Correlation Coefficient (MCC) for detecting SATD comments using GPT and Gemini across 0-, 2-, and 4-shot settings.}
\label{table:proprietary_llm_detection_result}
% \resizebox{\textwidth}{!}{
\begin{adjustbox}{max width=\textwidth, max height=\textheight}
\begin{tabular}{cccccccccccccc}
\toprule
\multirow{2}{*}{\textbf{Dataset}} & \multirow{2}{*}{\textbf{Model}} 
& \multicolumn{4}{c}{\textbf{0-shot}} 
& \multicolumn{4}{c}{\textbf{2-shots}} 
& \multicolumn{4}{c}{\textbf{4-shots}} \\
\cmidrule{3-14}
 &  & \textbf{P} & \textbf{R} & \textbf{F1} & \textbf{MCC} 
 & \textbf{P} & \textbf{R} & \textbf{F1} & \textbf{MCC} 
 & \textbf{P} & \textbf{R} & \textbf{F1} & \textbf{MCC}  \\
\midrule

\multirow{5}{*}{\rotatebox[origin=c]{45}{Original}} & gpt-5-nano  & 0.12 & 0.93 & 0.20 & 0.31  & 0.18 & 0.91 & 0.30 & 0.39  & 0.19 & 0.91 & 0.31 & 0.40  \\
 & gpt-5-mini  & 0.23 & 0.95 & 0.36 & 0.45  & 0.25 & 0.93 & 0.39 & 0.47  & 0.26 & 0.96 & 0.40 & 0.48  \\
 & gpt-5  & 0.27 & 0.95 & 0.42 & 0.50  & \textbf{0.34} & 0.93 & \textbf{0.49} & \textbf{0.55}  & 0.32 & 0.94 & 0.47 & 0.54  \\
\cmidrule(lr){2-2} \cmidrule(lr){3-14}
 & gemini-2.5-flash  & 0.15 & \textbf{0.99} & 0.26 & 0.36  & 0.17 & \textbf{1.00} & 0.28 & 0.39  & 0.12 & \textbf{0.99} & 0.21 & 0.33  \\
 & gemini-2.0-flash  & 0.24 & 0.88 & 0.37 & 0.44  & 0.15 & 0.98 & 0.26 & 0.37  & 0.12 & \textbf{0.99} & 0.21 & 0.32  \\
\midrule
\multirow{5}{*}{\rotatebox[origin=c]{45}{Deduplicate}} & gpt-5-nano  & 0.11 & 0.92 & 0.20 & 0.30  & 0.19 & 0.91 & 0.31 & 0.39  & 0.17 & 0.89 & 0.29 & 0.38  \\
 & gpt-5-mini  & 0.22 & 0.94 & 0.35 & 0.44  & 0.24 & 0.94 & 0.39 & 0.47  & 0.24 & 0.96 & 0.38 & 0.46  \\
 & gpt-5  & 0.27 & 0.96 & 0.42 & 0.49  & \textbf{0.31} & 0.92 & \textbf{0.46} & \textbf{0.52}  & 0.30 & 0.93 & 0.45 & \textbf{0.51}  \\
\cmidrule(lr){2-2} \cmidrule(lr){3-14}
 & gemini-2.5-flash  & 0.15 & \textbf{0.99} & 0.26 & 0.36  & 0.17 & \textbf{1.00} & 0.29 & 0.39  & 0.12 & \textbf{0.99} & 0.21 & 0.31  \\
 & gemini-2.0-flash  & 0.24 & 0.88 & 0.38 & 0.45  & 0.15 & 0.98 & 0.27 & 0.37  & 0.11 & \textbf{0.99} & 0.20 & 0.31  \\

\bottomrule

\end{tabular}
% }
\end{adjustbox}
\end{table*}


\textbf{Results.} The experimental results are summarized in Table~\ref{table:proprietary_llm_detection_result}. Proprietary models achieved consistently high recall, frequently exceeding 0.90 and occasionally reaching 1.0 (e.g., Gemini-2.5-Flash under 2-shot). However, this high recall was accompanied by very low precision, particularly for Gemini models, where precision often ranged between 0.11 and 0.24. Among GPT models, the standard GPT-5 demonstrated the most balanced performance, attaining the highest F1-scores under 2-shot settings on both the original (0.49) and deduplicated (0.46) datasets. The GPT-5 mini and nano variants performed moderately well, outperforming Gemini in terms of precision and F1, but lagging behind GPT-5 standard. Increasing the number of shots did not consistently improve performance, as was seen with the Flan-T5 models as well.

Compared to the Flan-T5 series (RQ3), proprietary models exhibit distinct performance patterns. GPT-5 models maintain high recall, unlike T5-XL, which achieves strong precision primarily with MAT or Jitterbug prompts but suffers from limited recall. GPT-5 provides more stable SATD detection across shot settings compared to the Flan-T5 series. Nonetheless, Flan-T5-XL with MAT prompts achieved the highest F1-score (0.68).

% \textcolor{red}{We were surprised by the overall poor performance of the LLMs, particularly the proprietary ones. This unexpected result prompted a closer examination, although a comprehensive analysis lies beyond the scope of this work. Nonetheless, even our limited investigation yielded several interesting observations. For instance, GPT-5 classified the comment \textit{“// implement the fake service”} as SATD, even though it merely describes the implementation of an anonymous service in the subsequent line. The model misinterpreted it as a placeholder for unfinished work. Such errors arise because LLMs do not rely solely on syntactic cues (e.g., keywords like TODO or FIXME) but attempt to infer semantic meaning. Even neutral or descriptive comments that express intentions, suggestions, or imperatives are frequently misclassified, inflating recall while substantially reducing precision.}

\begin{tcolorbox}
\textbf{\underline{RQ4 summary}:}
In general, proprietary LLMs, particularly GPT-5, outperform open-source models in recall. However, their much lower precision indicates a tendency to over-predict SATD instances. These results suggest that while proprietary models may offer stronger contextual reasoning to exhibit high recall, significant improvements are needed to address their limitations with precision.
\end{tcolorbox}
